\documentclass[10pt,a4paper]{amsart}
\usepackage[margin=19mm]{geometry}
\usepackage{amsmath,amssymb,mathtools}
\usepackage[scr=rsfs,cal=euler]{mathalfa}
\usepackage{xfrac}
\usepackage[unicode,hidelinks,bookmarks=false]{hyperref}
\newcommand{\ie}{i.e.\ }
\newcommand{\cf}{cf.\ }
\newcommand{\resp}{respectively\ }
\newcommand{\Subsection}{\subsection}
\providecommand{\nobreakdash}{\nobreak\hbox{-}}
\setlength{\emergencystretch}{2em}
\multlinegap0pt
\usepackage{bm}
\usepackage{bbm}
\usepackage{dsfont}
\usepackage{xspace}
\usepackage{cancel}
\usepackage{mathtools}
\usepackage{amsthm}
\makeatletter
\@ifundefined{lemma}{\newtheorem{lemma}{Lemma}}{}
\@ifundefined{remark}{\newtheorem{remark}[lemma]{Remark}}{}
\@ifundefined{definition}{\newtheorem{definition}[lemma]{Definition}}{}
\makeatother
\title[Notes on Laver Tables]{Notes on Laver Tables}
\author{Renrui Qi}
\date{Source: arXiv:2501.06733v5 (CC BY 4.0)}
\begin{document}
\maketitle

\noindent\emph{Source and licence.} Notes on Laver Tables. Version arXiv:2501.06733v5 (CC BY 4.0), distributed under the Creative Commons Attribution 4.0 International licence, as recorded in the arXiv licence metadata for this exact version and shown on the abstract page https://arxiv.org/abs/2501.06733v5. Full rights notice: licenses/arxiv-2501.06733-CC-BY-4.0.txt.\par\smallskip

\noindent\textit{Editorial scope.} Counted unit: the Lemma whose source label is \texttt{lem:palpha-successor-limit} in the article (source0.tex lines 1718--1724, source label \texttt{lem:palpha-successor-limit}), delivered together with the article's own complete original proof of it (source0.tex lines 1726--1843, one \texttt{proof} environment of 118 lines). No mathematics is written, repaired or re-derived: statement and proof are the article's own text line for line. Required context retained uncounted: rmk:applying\_a\_row\_of\_length\_3 (lines 737--739); def:limit\_and\_successor\_types (lines 778--789); def:operation\_E (lines 796--803); def:alpha-n (lines 1566--1574); lem:ps-basic (lines 1589--1598); eq:succ-tail (lines 1670--1672). Every definition, display and label of those lines is retained unchanged. Omitted from this package and not redistributed: 1--736, 740--777, 790--795, 804--1565, 1575--1588, 1599--1669, 1673--1717, 1725--1725, 1844--2401. Nothing inside a retained excerpt is omitted or altered: every retained body is delivered exactly as the article's lines are written, so no omission receipt of any kind accompanies this package. Citations: the retained excerpts carry no citation command at all, so no bibliography entry of the article is transcribed and no reference is redistributed. Reference adaptation: none was needed and none was made; every reference inside the delivered unit resolves inside it, so no unresolved reference or citation is produced. Printed slips kept literal: no printed word, symbol, hypothesis or number of the counted statement or of its proof was altered. Layout and class: the article's own class is \texttt{article} with packages including amsmath, amssymb, amsthm, xparse, tikz, listings, xcolor, caption, soul, hyperref; the wrapper is the lane's pinned \texttt{amsart} editorial wrapper at 10pt on A4, which restates the theorem-like environments the retained text needs and adds the article's own macro definitions that the retained text needs (none), with extra wrapper packages (bm, bbm, dsfont, xspace, cancel, mathtools). The delivered numbering is the unit's own. Assets: no retained span contains a figure environment, a graphics-inclusion command, a PDF-page inclusion, a table or a TikZ picture; no image file is copied into this package, the package declares no asset dependency and no image file is redistributed with it. Licence: version arXiv:2501.06733v5 of this article is distributed under the Creative Commons Attribution 4.0 International licence, as recorded in the arXiv licence metadata for this exact version and shown on the abstract page https://arxiv.org/abs/2501.06733v5. A full rights notice is delivered at licenses/arxiv-2501.06733-CC-BY-4.0.txt. Characters: every retained excerpt is pure ASCII, so the delivered engine, XeLaTeX with its default Latin Modern text font, reports no missing character.
\begin{remark} \label{rmk:applying_a_row_of_length_3}
  If the last row $s_n$ of~$p$ has length~3, and its first entry $s_{n,1}$ is nonzero, then $p$ is copyable, as $a = s_{n,1}$ (the step length~$\ell_n$ is necessarily~1). Since $s_{n,2} = n$ in this case, the rows to be copied would be $a,a+1,\dots, n-1$. The effect on each row would be to preserve the entries smaller than~$a$, and increse the other entries by $n-a$. 
\end{remark}

\begin{definition} \label{def:limit_and_successor_types} 
Let~$p$ be a blp; let $n = p.n$ be the length of~$p$.
\begin{enumerate}
  \item We say that $p$ is of \emph{limit type} if the first three entries of the last row $s_n$ of~$p$ are $0,1,2$, $|s_n|=6$, and the step length $\ell_n$ of the last row is~3.

  \item We say that $p$ is of \emph{successor type} if the first three entries of $s_n$ are $0,1,2$, and $|s_n|=5$ (so $s_n = (0,1,2,n,n+1)$, and $\ell_n=2$).

  \item The \emph{zero blp} is the unique blp of length~2. 

  \item If~$p$ is neither of limit type nor of successor type, and is not the zero blp, then we say that~$p$ is of \emph{transient type}.
\end{enumerate}
\end{definition}

\begin{definition} \label{def:operation_E}
  Suppose that~$p$ is a blp of either limit or successor type, of length~$n$.  Let $a=s_{n,-3}$. We define $p.E(m)$ by recursion on~$m$. 
  \begin{itemize}
    \item $p.E(0) = p.del$.

    \item Suppose that $p.E(m)$ has been defined, and has length~$n$. Let $q_{m+1}$ be the blp obtained from $p.E(m)$ by adding one row, $(a,n'+1,n'+2)$ where $n' = p.E(m).n$ is the length of~$p.E(m)$; the step length is necessarily~1. We let $p.E(m+1) = q_{m+1}.Copied$.\footnote{Note that by Remark~\ref{rmk:applying_a_row_of_length_3}, $p.E(m)$ is always well-defined.}
  \end{itemize}
\end{definition}

\begin{definition}\label{def:alpha-n}
Every nonzero $\alpha<\varepsilon_0$ admits a unique decomposition $\alpha=\beta+\omega^\gamma$ with $\gamma\ge 0$ and $\beta<\alpha$.
Define $\alpha[n]$ recursively by:
\begin{itemize}
  \item $\bigl(\omega^{\gamma+1}\bigr)[n]\ :=\ \omega^\gamma\cdot n$;
  \item if $\gamma$ is limit, $\bigl(\omega^\gamma\bigr)[n]\ :=\ \omega^{\gamma[n]}$;
  \item $(\beta+\omega^\gamma)[n]\ :=\ \beta+\bigl(\omega^\gamma\bigr)[n]$.
\end{itemize}
\end{definition}

\begin{lemma}\label{lem:ps-basic}
Let $\alpha<\varepsilon_0$, and write $ps(\alpha)=(u_1,\ldots,u_k)$.
Then:
\begin{enumerate}
  \item $ps(\alpha)=()$ iff $\alpha=0$;
  \item for all $i$, $0\le u_i<i$;
  \item $u_k=0$ iff $\alpha$ is a successor, in which case $ps(\alpha-1)=(u_1,\ldots,u_{k-1})$;
  \item if $\alpha$ is limit, then for all $n\ge 1$, $ps\bigl(\alpha[n]\bigr)$ is obtained from $ps(\alpha)$ by replacing the final block according to Definition~\ref{def:alpha-n} (hence $ps\bigl(\alpha[n]\bigr)$ is an initial segment of $ps\bigl(\alpha[n+1]\bigr)$).
\end{enumerate}
\end{lemma}

\begin{multline}\label{eq:succ-tail}
ps(\alpha[n])\ =\ A\ \frown\ \bigl(\ 0\ \frown\ ps(\delta)^{+(a)}\ \bigr)\ \frown\ \bigl(\ 0\ \frown\ ps(\delta)^{+\,(a+(1+|ps(\delta)|))}\ \bigr)\ \frown\ \\ \cdots\ \frown\  \bigl(\ 0\ \frown\ ps(\delta)^{+\,(a+(n-1)(1+|ps(\delta)|))}\ \bigr).
\end{multline}

\begin{lemma}\label{lem:palpha-successor-limit}
For all $\alpha<\varepsilon_0$ and $m\ge 0$:
\begin{enumerate}
  \item $p_{\alpha+1}$ is of successor type and $(p_{\alpha+1}).{\mathrm{del}}=p_\alpha$;
  \item if $\alpha$ is limit, then $p_\alpha.E(m)=p_{\alpha[2^m]}$.
\end{enumerate}
\end{lemma}

\begin{proof}
Recall that $p_\alpha$ is defined from $ps(\alpha)=(t_1,\dots,t_k)$ by taking
$p_\alpha.n=k+2$ and, for $1\le i\le k$,
\[
t_i=0 \ \Rightarrow\ p_\alpha.s_{i+2}=(0,1,2,i+2,i+3),\quad
t_i>0 \ \Rightarrow\ p_\alpha.s_{i+2}=(0,1,2,t_i+2,i+2,i+3),
\]
with step lengths $2$ and $3$ respectively. The first two rows are fixed
$(0,1,2)$ and $(0,1,2,3)$.

\smallskip
\noindent\textbf{(i) } $p_0$ is the zero blp; for every $\alpha$, $p_{\alpha+1}$ is of successor type and $(p_{\alpha+1})\!.del=p_\alpha$.

Since $ps(0)=()$ (Lemma~\ref{lem:ps-basic}(1)), $p_0$ has exactly the two fixed rows, so it is the zero blp.
For $\alpha+1$, Lemma~\ref{lem:ps-basic}(3) gives
\(
ps(\alpha+1)=ps(\alpha)\frown 0
\)
and $ps(\alpha)$ is exactly the prefix obtained by deleting the final entry.
Therefore $p_{\alpha+1}$ is obtained from $p_\alpha$ by appending one new row
corresponding to $t_{k+1}=0$, namely $(0,1,2,(k+2),(k+3))$ with step length $2$.
By Definition~\ref{def:limit_and_successor_types}(2) this makes $p_{\alpha+1}$ of
successor type, and deleting that last row gives back $p_\alpha$, i.e.
$(p_{\alpha+1})\!.del=p_\alpha$.

\smallskip
\noindent\textbf{(ii) } If $\alpha$ is limit, then for all $m\in\mathbb{N}$,
\(
p_\alpha.E(m)=p_{\alpha[2^m]}.
\)

Fix a limit $\alpha$. Write $\alpha=\beta+\omega^\gamma$ with $\gamma>0$ and set
$A:=ps(\beta)$, $a:=|A|+1$. By Lemma~\ref{lem:ps-basic}(2), every entry of $ps(\alpha)$ is
$<\,$its position, so $p_\alpha$ is well-defined. By the same lemma and the
definition of successor/limit types, the last row $s_{n}$ of $p_\alpha$ begins with
$(0,1,2,\dots)$ and:
\begin{itemize}
  \item if $\gamma=\delta+1$ (successor tail), then $ps(\alpha)=A\frown 0\frown ps(\delta)^{+(a)}\frown a$ and $p_\alpha$ is of \emph{limit} type with $s_{n,-3}=a$;
  \item if $\gamma$ is limit, then $ps(\alpha)=A\frown 0\frown ps(\gamma)^{+(a)}$ and $p_\alpha$ is of \emph{limit} type with $s_{n,-3}=a$ as well (here the final block has length $\ge1$ and ends at position $n$, so its third-from-last entry exists and equals $a$).
\end{itemize}
Thus in both subcases the parameter $a$ of Definition~\ref{def:operation_E} equals
$s_{n,-3}$.

We prove by induction on $m$ that $p_\alpha.E(m)=p_{\alpha[2^m]}$.

\emph{Base $m=0$.}
By Definition~\ref{def:operation_E}, $p_\alpha.E(0)=p_\alpha.del$.
By Lemma~\ref{lem:ps-basic}(4), when $\alpha$ is limit one has
\(
ps(\alpha[1])
\)
obtained from $ps(\alpha)$ by replacing its last block with an appropriate initial
segment; in particular $ps(\alpha[1])$ is exactly $ps(\alpha)$ with the final entry
removed in the successor-tail case, and the corresponding ``one-step'' truncation in
the limit-tail case. In both subcases this is precisely the effect of $p\mapsto p.del$
on the sequence of rows. Hence $p_\alpha.E(0)=p_\alpha.del=p_{\alpha[1]}$.

\emph{Step $m\to m+1$.}
Assume $p_\alpha.E(m)=p_{\alpha[2^m]}$. Write $q:=p_{\alpha[2^m]}$ and let
$q.n$ be its length. By Definition~\ref{def:operation_E}, to compute $p_\alpha.E(m+1)$
we start from $q$, append one row $(a,q.n+1,q.n+2)$ (with the same $a=s_{q.n,-3}$ as
above), obtaining $q_{m+1}$, and then apply the \textit{Copied} operation to $q_{m+1}$
with parameter $a$.

We analyze the effect on the tail determined by the Cantor-normal-form tail of
$\alpha[2^m]$ and compare with Lemma~\ref{lem:ps-basic}(4).

\underline{Case 1: } $\gamma=\delta+1$ (successor tail).
By Lemma~\ref{lem:ps-basic}(4) (equation~\eqref{eq:succ-tail} in its proof), for all $n\ge1$
\begin{multline*}
  ps\!\bigl((\beta+\omega^{\delta+1})[n]\bigr)
= \\ A \frown \bigl(0\frown ps(\delta)^{+(a)}\bigr)
      \frown \bigl(0\frown ps(\delta)^{+(a+(1+|ps(\delta)|))}\bigr)
      \frown \cdots \frown \\ 
      \bigl(0\frown ps(\delta)^{+(a+(n-1)(1+|ps(\delta)|))}\bigr).
\end{multline*}
In particular,
\[
ps\!\bigl(\alpha[2^m]\bigr)=A\frown \underbrace{B \frown B \frown\cdots \frown B}_{2^m\ \text{blocks}},
\quad\text{where }B:=\bigl(0\frown ps(\delta)^{+(a)}\bigr),
\]
and hence $ps\!\bigl(\alpha[2^{m+1}]\bigr)$ is obtained by \emph{doubling} the number
of copies of $B$ from $2^m$ to $2^{m+1}$.

Now observe what $E$ does on $q=p_{\alpha[2^m]}$:
appending the row $(a,q.n+1,q.n+2)$ and applying \textit{Copied} with the same $a$
copies \emph{all rows indexed $\ge a$} and shifts their entries $\ge a$ by a fixed
stride (Remark~\ref{rmk:applying_a_row_of_length_3}). In the canonical $p_\bullet$
encoding, the rows $\ge a$ are exactly the current $2^m$ copies of the block $B$.
Therefore, $q\mapsto q_{m+1}\mapsto q_{m+1}.\mathrm{Copied}$ duplicates those $2^m$
blocks into $2^{m+1}$ blocks, without changing the prefix determined by $A$.
Thus the resulting blp has pattern sequence $ps\!\bigl(\alpha[2^{m+1}]\bigr)$, i.e.
$p_\alpha.E(m+1)=p_{\alpha[2^{m+1}]}$.

\underline{Case 2: } $\gamma$ limit.
By Lemma~\ref{lem:ps-basic}(4),
\[
ps\!\bigl(\alpha[n]\bigr)=A\frown 0 \frown ps\!\bigl(\gamma[n]\bigr)^{+(a)},
\]
and $ps(\gamma[n])$ is an initial segment of $ps(\gamma[n+1])$ for all $n$.
Hence going from $n=2^m$ to $n=2^{m+1}$ \emph{replaces} the last block
$ps(\gamma[2^m])^{+(a)}$ by a longer block $ps(\gamma[2^{m+1}])^{+(a)}$.

On the $p_\bullet$ side, $q=p_{\alpha[2^m]}$ has the form
\(
A\frown 0 \frown C^{+(a)}
\)
with $C=ps(\gamma[2^m])$. Appending $(a,q.n+1,q.n+2)$ and applying
\textit{Copied} with parameter $a$ copies exactly the rows in the last block and
extends it \emph{by the next chunk determined by the growth of $ps(\gamma[n])$}.
By Lemma~\ref{lem:ps-basic}(4) for $\gamma$, after one such copy the last block
becomes $ps(\gamma[2^{m+1}])^{+(a)}$. The prefix determined by $A\frown 0$ is
unchanged. Therefore the resulting blp has pattern sequence
$ps\!\bigl(\alpha[2^{m+1}]\bigr)$, i.e.\ $p_\alpha.E(m+1)=p_{\alpha[2^{m+1}]}$.

\smallskip
This completes the induction on $m$ and proves (ii).
\end{proof}

\end{document}
